\documentclass[10pt]{article}
\usepackage[french]{babel}
\usepackage[utf8]{inputenc}
\usepackage[T1]{fontenc}
\usepackage{amsmath}
\usepackage{amsfonts}
\usepackage{amssymb}
\usepackage[version=4]{mhchem}
\usepackage{stmaryrd}

\begin{document}
\section*{Analyse Numérique - TD n \({ }^{\circ} 4\) 20-03-2025}
\section*{Exercice 1:}
Soit l'équation de la chaleur suivante :

\[
\left.\frac{\partial \theta}{\partial t}=K \frac{\partial^{2} \theta}{\partial x^{2}} \quad, \quad x \in\right] 0, L[, \quad K \text { est une constante }
\]

Conditions initiales :

\[
t=0 ; \forall x \in] 0, L\left[; \quad \theta(x, 0)=\theta_{0}=\alpha x^{2}+1 ; \quad \alpha\right. \text { est une constante }
\]

Conditions aux limites \(\forall t\) :

\[
x=0 ; \frac{\partial \theta}{\partial x}=0 ; \quad x=L ; \quad \frac{\partial \theta}{\partial x}=2 \alpha L
\]

\begin{enumerate}
  \item Vérifier que le problème est bien posé avec les conditions initiales et aux limites proposées
  \item Ecrire l'équation discrétisée en associant correctement à chacun des termes de l'équation, les schémas numériques ci-dessous :
\end{enumerate}

\begin{itemize}
  \item schéma en temps d'Adams implicite : \(\theta^{n+1}=\theta^{n}+\Delta t\left[\frac{3}{4} L^{n+1}+\frac{1}{4} L^{n-1}\right]\)
  \item schéma en temps Euler implicite : \(\theta^{n+1}=\theta^{n}+\Delta t L^{n+1}\)
  \item schéma aux différences-finies centrée pour la dérivée seconde d'espace
\end{itemize}

\[
\left(\frac{\partial^{2} \theta}{\partial x^{2}}\right)_{i}=\frac{\theta_{i+1}-2 \theta_{i}+\theta_{i-1}}{\Delta x^{2}}
\]

\begin{itemize}
  \item schémas aux différences-finies décentrée pour la dérivée première d'espace
\end{itemize}

\[
\left(\frac{\partial \theta}{\partial x}\right)_{i}=\frac{\theta_{i+1}-\theta_{i}}{\Delta x} \quad \text { et } \quad\left(\frac{\partial \theta}{\partial x}\right)_{i}=\frac{\theta_{i}-\theta_{i-1}}{\Delta x}
\]

\(\Delta \mathrm{t}, \Delta \mathrm{x}\) sont respectivement le pas de temps et le pas d'espace. Ils sont constants.\\
3. Ecrire le problème matriciel obtenu pour une discrétisation spatiale de 6 points (y compris les bornes du domaine de définition) régulièrement espacés.\\
4. Peut-on résoudre le problème matriciel par une méthode itérative (Jacobi ou Gauss-Seidel) ?\\
5. On modifie les conditions initiales et aux limites :

Conditions initiales :

\[
t=0 ; \forall x \in] 0, L\left[; \theta(x, 0)=\theta_{0}=1+3 x^{2}\right.
\]

\section*{Conditions aux limites \(\forall t\) :}
\[
x=0 ; \theta=1 ; \quad x=L ; \quad \frac{\partial \theta}{\partial x}=6 L
\]

Ecrire de nouveau le problème matriciel pour une discrétisation spatiale de 6 points avec ces nouvelles conditions aux limites.

\section*{Exercice 2:}
L'équation de Burgers s'écrit :

\[
\left.\frac{\partial u}{\partial t}+u \frac{\partial u}{\partial x}=v \frac{\partial^{2} u}{\partial x^{2}} \quad x \in\right] 0, L[
\]

avec des conditions aux limites et des conditions initiales.\\
Cette équation décrit l'évolution d'un champ de vitesse \(u(x, t)\) sous l'influence des deux termes :

\begin{itemize}
  \item Un terme d'advection \(u \frac{\partial u}{\partial x}\)
  \item Un terme de dissipation linéaire \(v \frac{\partial^{2} u}{\partial x^{2}}\), où \(v\) est une constante.
\end{itemize}

\begin{enumerate}
  \item Ecrire l'équation discrétisée en associant correctement à chacun des termes de l'équation, les schémas numériques ci-dessous
\end{enumerate}

\begin{itemize}
  \item schéma en temps Euler explicite (pour la partie non linéaire)
\end{itemize}

\[
u^{n+1}=u^{n}+H^{n} \Delta t
\]

\begin{itemize}
  \item schéma en temps Euler implicite (pour la partie linéaire)
\end{itemize}

\[
u^{n+1}=u^{n}+H^{n+1} \Delta t
\]

\begin{itemize}
  \item schémas des différences-finies centrées pour les deux dérivées d'espace
\end{itemize}

\[
\left(\frac{\partial u}{\partial x}\right)_{i}=\frac{u_{i+1}-u_{i-1}}{2 \Delta x}, \quad\left(\frac{\partial^{2} u}{\partial x^{2}}\right)_{i}=\frac{u_{i+1}-2 u_{i}+u_{i-1}}{\Delta x^{2}}
\]

\(\Delta \mathrm{t}, \Delta \mathrm{x}\) sont respectivement le pas de temps et le pas d'espace. Ils sont constants.\\
2. A quelle condition le problème est-il bien posé ? On choisit des conditions aux limites de Dirichlet \(u=u_{0}=\) constante.\\
3. Ecrire le problème matriciel ainsi obtenu pour une discrétisation spatiale de 6 points (y compris les bornes du domaine de définition) régulièrement espacés. Suggérer une méthode de résolution et justifier la.


\end{document}